% ============================================================
% conclusion.tex — Conclusion and Future Directions
% ============================================================

This chapter synthesizes the central findings of the thesis,
identifies the most productive avenues for extending this
work, and reflects on the broader implications of the
competing-risks approach for empirical legal scholarship.
Section~7.1 summarizes the four core substantive
contributions,
Section~7.2 proposes specific research directions motivated
by the limitations identified in
Chapter~\ref{ch:discussion}, and Section~7.3 offers a
closing reflection.

\section{Summary of Contributions}

This thesis asked what determines the timing and type of
resolution in federal securities class actions. The
characterization of that answer depends critically on the
statistical framework applied. By applying a full
competing-risks survival
framework to 12,968 cases from the FJC Integrated Database
spanning 1990--2024, we have shown that the standard
single-outcome approach systematically distorts the picture
of securities litigation dynamics.

Building on the methodological framework developed in
Chapters~\ref{ch:methodology}--\ref{ch:data} and the
analytical results presented in
Chapters~\ref{ch:results}--\ref{ch:discussion}, four
substantive contributions emerge.

First, the PSLRA's association with case outcomes is
time-varying rather than uniform. The dismissal association is
concentrated early and fades over litigation duration,
whereas the settlement-side piecewise model remains below
unity but sits within a much less stable broader settlement
record. Taken together, the most defensible substantive claim
is an early dismissal filter, not a uniform shift in all
aspects of litigation dynamics.

Second, the IPTW composition-adjusted sensitivity analysis
(Section~\ref{sec:results_iptw}) shows that conditioning on
observable case-mix shifts leaves the dismissal association
elevated and materially attenuates the settlement
association. Those contrasts are informative, but they are
sensitivity comparisons rather than a formal decomposition of
the raw PSLRA gap, especially because some weighting
covariates may themselves lie on post-PSLRA pathways. The
settlement association in the extended fixed-effects model
also does not survive the cluster-robust specification
($p = 0.158$), and the settlement result weakens in the two
largest circuit-specific sub-samples. The IPTW estimates are
composition-adjusted rather than causally identified, but the
broader triangulation provides a principled sensitivity check
unavailable in prior work.

Third, circuit geography emerges as a dominant structural
predictor of case trajectories. Relative to the Second
Circuit, every other major circuit is more settlement-prone,
and the frailty results suggest that settlement outcomes are
more sensitive than dismissal outcomes to local judicial
culture and bar practices.

Fourth, the MDL sign-reversal finding
(Section~\ref{sec:results_characteristics}) illustrates why
competing-risks analysis is especially valuable here. MDL
cases resolve more slowly at the cause-specific level but
have higher cumulative settlement probability once the
competing dismissal pathway is modeled directly.

Throughout, we have treated the 1992 placebo failure, the
post-treatment-bias concern in the IPTW covariates, and the
settlement association's fragility in circuit-specific
subsamples as binding interpretive limits rather than as
afterthoughts. Those boundaries shape the contribution, but
they are not themselves a separate substantive finding.

\section{Future Directions}

Among these extensions, the most tractable near-term steps
are district-level disaggregation within the existing IDB and
narrow-window re-estimation using the current covariate set;
CourtListener, Stanford, and CRSP merges are more ambitious
data-integration projects.

\subsection{Unmeasured Confounders and Data Enrichment}

The most binding limitation of this thesis is the absence of
case-specific merit signals from the FJC's administrative
data. Three categories of unmeasured confounders warrant
investigation in future work.

\paragraph{Judicial ideology and case-management norms.}
The frailty models (Section~\ref{sec:frailty}) detect
suggestive unobserved settlement heterogeneity at the
circuit level ($\hat{\theta} = 0.130$), but cannot identify its source.
Judicial ideology scores---such as the Judicial Common Space
measures \citep{epstein2007}---combined with judge
identifiers recoverable from CourtListener docket records,
would allow extension of the current circuit-level frailty
to judge-level random effects. This would permit direct
estimation of whether individual judicial approaches to case
management, discovery scheduling, and motion practice drive
the geographic patterns we observe.

\paragraph{Plaintiff firm sophistication.}
\citep{wang2022} documents that law firm identity is a
strong predictor of litigation outcomes, yet the IDB lacks
attorney identifiers. CourtListener integration would
provide title-based class-action validation, lead-counsel
identification, and docket-activity proxies for litigation
complexity. These variables could be incorporated as
additional covariates or as stratification factors in the
competing-risks models, sharpening interpretation of the
PSLRA's filtering association.

\paragraph{Economic conditions and market context.}
Securities fraud allegations are often triggered by sharp
stock-price declines, and the business cycle may influence
both the volume and the composition of filings. Merging with
CRSP market data would enable examination of whether
economic conditions at the time of filing moderate the PSLRA
association, and whether the post-PSLRA compositional shift
that our IPTW analysis adjusts for is partly driven by
market-cycle effects not captured by the current covariates.

\subsection{The Second Circuit Pattern}

The most striking empirical finding in our geographic
analysis is not that the Second Circuit has the single
highest dismissal CIF, but that it combines the lowest
settlement baseline with persistently high dismissal
incidence. Every other major circuit's extended settlement
hazard is 1.6 to 3.5 times that of the Second Circuit, while
the Fifth Circuit slightly exceeds the Second in cumulative
dismissal incidence. We have offered potential
mechanisms---elite defense counsel, judicial expertise,
demanding class-certification standards---but cannot
disentangle them with the current data.

A dedicated investigation of this pattern would be valuable.
First, disaggregating the Second Circuit into its
constituent districts (SDNY, EDNY, NDNY, WDNY, and the
District of Vermont) would reveal whether the pattern is
driven primarily by the Southern District of New York or is
circuit-wide. Second, comparing pre- and post-PSLRA trends
within the Second Circuit alone could clarify whether the
circuit's distinctive environment amplified or moderated the
reform's association with outcomes. Third, merging with the
Stanford Securities Class Action Clearinghouse would provide
settlement amount data that could test whether the Second
Circuit's low settlement hazard reflects not fewer
settlements but slower ones---cases that take longer to
settle but ultimately resolve for larger amounts.
The district-level disaggregation is the most immediately
tractable of these extensions because it can be executed
within the existing IDB.

\subsection{Methodological Extensions}

Several methodological extensions would address limitations
identified in this thesis. Full time-varying coefficient
models for circuit, origin, MDL, and statutory basis are
needed to address the documented proportional hazards
assumption violations that affect all constant hazard-ratio
estimates. A piecewise IPTW model---decomposing the
composition-adjusted PSLRA association across the same three
time periods used in the unweighted analysis---would test
whether the time-varying signature survives composition
adjustment. A narrow-window IPTW analysis restricted to
cases filed 1993--1998 would provide the strongest
quasi-experimental test of the PSLRA's association, though
the smaller sample may limit statistical power.

Settlement amount analysis, made possible by merging with the
Stanford Securities Class Action Clearinghouse, would extend
the competing-risks duration framework into a joint model of
timing, outcome type, and economic magnitude---providing a
more complete picture of the PSLRA's full association with
the securities litigation system.

\subsection{Structural Breaks and Emerging Dynamics}

Testing for post-2018 structural breaks---particularly the
Supreme Court's 2018 decision in \textit{Cyan, Inc.\ v.\
Beaver County Employees Retirement Fund}\footnote{140 S. Ct.\
1061 (2018).}, which expanded
state
court jurisdiction over Securities Act claims---is a natural
extension of the piecewise hazard approach developed here.
If \textit{Cyan} redirected Section~11 claims from federal
to state courts, the composition of the federal securities
docket may have shifted in ways that the competing-risks
framework is well positioned to detect. Similarly, the rapid
growth of SPAC- and cryptocurrency-related litigation since
2020 may introduce a new class of cases with distinctive
duration and outcome patterns that warrant separate
analysis within the competing-risks framework.

\section{Closing Reflection}

The central methodological lesson of this thesis is that the
competing-risks framework materially changes the empirical
picture of securities litigation dynamics. A standard
single-outcome survival analysis
of the same data would have found a significant PSLRA
association with dismissal timing, a significant settlement
suppression, and meaningful circuit variation---but it would
have missed the MDL sign reversal, mischaracterized the
PSLRA's settlement association, and provided no tools for
separating the reform's association from concurrent
compositional shifts.
The competing-risks framework does not merely refine the
estimates; it reveals structural features of the litigation
landscape that are invisible to the standard approach.

More broadly, this thesis demonstrates that empirical legal
scholarship benefits from the same methodological rigor that
has transformed biomedical and actuarial research. The
IDB represents one
of the most comprehensive longitudinal records of federal
civil litigation ever assembled, yet it has been
underutilized by researchers with formal training in survival
analysis and competing-risks methods. We hope that the
framework developed here---nonparametric CIF estimation,
cause-specific and subdistribution models, IPTW composition
adjustment, and frailty sensitivity analysis---serves as a
template for future empirical work on federal litigation, in
securities law and beyond. The same competing-risks
structure---multiple mutually exclusive outcomes, large
administrative datasets, and the need to distinguish
time-to-event from event-type---applies directly to
bankruptcy proceedings, employment discrimination cases, and
patent disputes, where the analytical challenges of the
present thesis recur in different legal guise.

The federal securities docket does not resolve neatly. Cases
settle, are dismissed, linger in MDL consolidation, or exit
for reasons the data cannot explain. Any framework that
pretends otherwise---by treating one outcome as the only
outcome, or by averaging over a time-varying process---will
mischaracterize the system it purports to study. A
competing-risks approach takes that complexity seriously, and
in doing so, comes closer to the truth.
